\documentclass[letterpaper, 10 pt, journal, twoside]{IEEEtran}

\IEEEoverridecommandlockouts %


\markboth{IEEE Robotics and Automation Letters. Preprint Version. Accepted June, 2020}
{Lusk \MakeLowercase{\textit{et al.}}: A Distributed Pipeline for Scalable, Deconflicted Formation Flying} 



%% the rest of your preamble here
\newcommand\blfootnote[1]{%
  \begingroup
  \renewcommand\thefootnote{}\footnote{#1}%
  \addtocounter{footnote}{-1}%
  \endgroup
}







\newcommand\eva[1]{\textcolor{pink}{Eva: #1}}
\newcommand\toptile[1]{top_{#1}}
\newcommand\bottomtile[1]{bot_{#1}}
\newcommand\lefttile[1]{lef_{#1}}
\newcommand\righttile[1]{rig_{#1}}
\newcommand{\coms}{\xspace\scalebox{0.5}{\begin{tikzpicture}\draw[dash pattern=on 3pt off 3pt, line width = 1mm ] (0,0) -- (0.8,0); \node at (0,-0.08) {}; \end{tikzpicture}}\xspace}
\newcommand{\moves}{\rightarrow}
\newcommand{\colorellipse}{orange!20!white}
\newcommand{\addCommunicationB}[1]
{
	\node[above right = 0.5mm and 4mm of #1, inner sep=0mm] (nb2) {B};
	\draw[tobehere,communication] (#1) -- (nb2);
}

\newcommand\setnodes{V}
\newcommand\basenode{B}
\newcommand{\sourcenode}{s}
\newcommand{\targetnode}{t}
\newcommand\USTCONN{USTCONN}
\newcommand\UCONN{UCONN}
\newcommand{\myundirectedgraphbounded}{\mathfrak G}
\newcommand{\myundirectedgraph}{\mathfrak G'}

\newcommand\pb[2]{${#1}_{#2}$~}

\newcommand\pbReach[1]{\pb{Reachability}{#1}}
\newcommand\pbBReach[1]{\pb{bReachability}{#1}}
\newcommand\pbCoverage[1]{\pb{Coverage}{#1}}
\newcommand\pbBCoverage[1]{\pb{bCoverage}{#1}}

\newcommand{\idDir}{dir}	% identifier for directed topological graph
\newcommand{\idNC}{nc}		% identifier for neighbor-communicable topological graph
\newcommand{\idSM}{sm}		% identifier for sight-moveable topological graph
\newcommand{\idCC}{cc}		% identifier for complete-communication topological graph
\newcommand{\idUnd}{und}		% identifier for complete-communication topological graph

\newcommand{\pbReachDir}{\pbReach{\idDir}}
\newcommand{\pbBReachDir}{\pbBReach{\idDir}}
\newcommand{\pbCoverageDir}{\pbCoverage{\idDir}}
\newcommand{\pbBCoverageDir}{\pbBCoverage{\idDir}}

\newcommand{\pbReachNC}{\pbReach{\idNC}}
\newcommand{\pbBReachNC}{\pbBReach{\idNC}}
\newcommand{\pbCoverageNC}{\pbCoverage{\idNC}}
\newcommand{\pbBCoverageNC}{\pbBCoverage{\idNC}}

\newcommand{\pbReachSM}{\pbReach{\idSM}}
\newcommand{\pbBReachSM}{\pbBReach{\idSM}}
\newcommand{\pbCoverageSM}{\pbCoverage{\idSM}}
\newcommand{\pbBCoverageSM}{\pbBCoverage{\idSM}}

\newcommand{\pbReachCC}{\pbReach{\idCC}}
\newcommand{\pbBReachCC}{\pbBReach{\idCC}}
\newcommand{\pbCoverageCC}{\pbCoverage{\idCC}}
\newcommand{\pbBCoverageCC}{\pbBCoverage{\idCC}}

\newcommand{\pbReachUnd}{\pbReach{\idUnd}}
\newcommand{\pbBReachUnd}{\pbBReach{\idUnd}}
\newcommand{\pbCoverageUnd}{\pbCoverage{\idUnd}}
\newcommand{\pbBCoverageUnd}{\pbBCoverage{\idUnd}}

\newtheorem{theorem}{Theorem}
\newtheorem{fact}[theorem]{Fact}
\newtheorem{definition}[theorem]{Definition}
\newtheorem{proposition}[theorem]{Proposition}
%\newtheorem{proof}{Proof}
\newtheorem{lemma}[theorem]{Lemma}

\tikzstyle{copyofG} = [fill=gray!20!white, draw=none, rounded corners=8pt]
\tikzstyle{communication} = [dash pattern=on 1pt off 1pt] %[line cap=rect, dot diameter=1pt, dot spacing=2pt, dots] %[dash pattern=on 2pt off 2pt on \the\pgflinewidth off 2pt]
\tikzstyle{tobehere} = [color=blue!50!white]
\newcommand\set[1]{\{#1\}}
\newcommand\suchthat{\mid}
\newcommand{\nbdrones}{n}

\newcommand{\booleanvariablegeneric}{{x}}
\newcommand{\clausegeneric}{\mathsf{c}}

\newcommand{\ok}{\includegraphics[height=2mm]{check.png}}
\newcommand{\imgDrone}{\includegraphics[height=2mm]{drone_black.png}}

\usetikzlibrary{decorations.pathmorphing}
\usetikzlibrary{decorations.pathreplacing}
\usetikzlibrary{arrows}
\usetikzlibrary{positioning}

\makeatletter
\tikzset{
	dot diameter/.store in=\dot@diameter,
	dot diameter=3pt,
	dot spacing/.store in=\dot@spacing,
	dot spacing=10pt,
	dots/.style={
		line width=\dot@diameter,
		line cap=round,
		dash pattern=on 0pt off \dot@spacing
	}
}
\makeatother

%%% Local Variables:
%%% mode: latex
%%% TeX-master: "main"
%%% End:

\pdfminorversion=4

\newcommand{\edittwo}[1]{\textcolor{black}{#1}}
\newcommand{\edit}[1]{\textcolor{black}{#1}}

\title{%
	 A Distributed Pipeline for \\ Scalable, Deconflicted Formation Flying
	}

\author{Parker C. Lusk, Xiaoyi Cai, Samir Wadhwania, Aleix Paris, Kaveh Fathian, Jonathan P. How%
    \thanks{Manuscript received: February 24, 2020; Revised: May 20, 2020; Accepted: June 1, 2020}%
	\thanks{
	    This paper was recommended for publication by N.Y.\ Chong upon evaluation of the Associate Editor and Reviewers' comments. %
	    Research supported in part by NASA Convergent Aeronautics Solutions project Design Environment for Novel Vertical Lift Vehicles (DELIVER),
		ARL DCIST under Cooperative Agreement Number W911NF-17-2-0181,
		and Boeing Research \& Technology.
		Computation support provided by Amazon Web Services.}
	\thanks{P.\ C.\ Lusk, X.\ Cai, S.\ Wadhwania, A.\ Paris,  K.\ Fathian and J.\ P.\ How are with the Department of Aeronautics and Astronautics, Massachusetts Institute of Technology.
	    {\{plusk, xyc, samirw, aleix, kavehf, jhow\}@mit.edu.}}
	\thanks{Digital Object Identifier (DOI): see top of this page.} %
}%

\begin{document}

\maketitle


\begin{abstract}
Reliance on external localization infrastructure and centralized coordination are main limiting factors for formation flying of vehicles in large numbers and in unprepared environments.
While solutions using onboard localization address the dependency on external infrastructure, the associated coordination strategies typically lack collision avoidance and scalability.
To address these shortcomings, we present a unified pipeline with onboard localization and a distributed, collision-free formation control strategy that scales to a large number of vehicles.
Since distributed collision avoidance strategies are known to result in gridlock, we also present a \edit{distributed} task assignment solution to deconflict vehicles.
We experimentally validate our pipeline in simulation and hardware.
The results show that our approach for solving the optimization problem associated with formation control gives solutions within seconds in cases where general purpose solvers fail due to high complexity.
In addition, our lightweight assignment strategy leads to successful and quicker formation convergence in
\SIrange[range-units=single,range-phrase=--]{96}{100}{\percent} of all trials,
whereas indefinite gridlocks occur without it for \SIrange[range-units=single,range-phrase=--]{33}{50}{\percent} of trials.
By enabling large-scale, deconflicted coordination, this pipeline should help pave the way for anytime, anywhere deployment of aerial swarms.
\end{abstract}


\begin{IEEEkeywords}
Swarms; Distributed Robot Systems; Multi-Robot Systems
\end{IEEEkeywords}
\vspace{-0.8em}


 \section*{Supplementary Material}


Video and open-source implementation available at \href{https://github.com/mit-acl/aclswarm}{\texttt{https://github.com/mit-acl/aclswarm}}.

\section{Introduction}\label{sec:intro}


\IEEEPARstart{T}{wo} %
main challenges in the deployment of large-scale swarms are the localization and coordination of vehicles.
Localization methods that rely on external infrastructure 
(e.g., GPS) 
are prone to systematic errors (e.g., multipath effect)
and may not always be available.
Coordination strategies that are centralized can deconflict motion plans to prevent collisions and gridlock, but introduce a single point of failure and are difficult to scale in swarm size due to communication bandwidth limitations.

This paper presents a unified formation flying pipeline for unmanned aerial vehicles (UAVs).
Our pipeline uses \textit{onboard} sensors for localization, which eliminate the need for external positioning systems, and \textit{distributed} techniques for coordination, which enable each vehicle to make decisions independently while communicating their state to a subset of the team.
For \textit{localization}, we use an off-the-shelf commercial visual inertial odometry (VIO) package \cite{VIO}
that fuses inertial measurement unit (IMU) and downward-facing monocular camera measurements to estimate changes in the vehicle pose.
\edit{For \textit{coordination}, we present distributed formation control and task assignment strategies that run onboard the vehicles, do not rely on a common reference frame, and use vehicle-to-vehicle communication.} 
Key features of our formation control strategy include scalability to a large number of vehicles and robustness to disturbances.
The latter is crucial for reaching the desired formations with sensing imperfections.
Our task assignment strategy uses an auction-based algorithm to guarantee conflict-free assignments.
This algorithm can deconflict vehicle gridlocks resulting from distributed collision avoidance (type 3 deadlock~\cite{Wang2017}) and is well-suited for vehicles with limited computational capability and low-bandwidth communication. 


\begin{figure}[t!]
	\begin{center}
		\includegraphics[trim =0mm 10mm 0mm 0mm, clip, width=\columnwidth]{Figs/slanted_plane.png}	
		\caption{
		Six multirotors in a slanted plane formation.
		Vehicles communicate with each other, make distributed decisions onboard, and use VIO for localization.}
		\label{fig:slantedplane}
	\end{center}
\end{figure}


\subsection{Contributions}

This research extends our previous work on UAV formations~\cite{Fathian2019} and presents a unified pipeline consisting of \textit{onboard localization} and \textit{distributed coordination}.
The three main contributions of this work are:
\begin{enumerate}
    \item \edit{scalable formulation of control design suitable for
    onboard sensing without a common reference frame;}
    \item algorithms for deconfliction via \edit{distributed} task assignment of vehicles to desired formation points;    
    \item simulation- and hardware-ready open-source pipeline.
\end{enumerate}
\edit{Our pipeline is tested in hardware with six multirotors (see Fig.~\ref{fig:slantedplane}), and 
to our knowledge is the first demonstration of formation flying that does not rely on external sensing, fiducial markers for localization, a common reference frame, or a centralized base station for coordination.}
The only requirements for the presented pipeline are that the vehicles can communicate, can find the transformation between their VIO start frames, and the environment is sufficiently textured---a standard assumption for VIO systems.
As such, this framework paves the way for future, real-world deployments of aerial vehicle swarms in large numbers and without requiring external localization infrastructure.


\begin{figure} [t!]
\centering
	\begin{subfigure}[b]{0.32\columnwidth}
	   %
	    \includegraphics[width=0.8\textwidth,left]{Figs/Frames2_full.pdf}
	    \caption{\scriptsize full alignment}
	    \label{fig:frame-a}
	\end{subfigure}
	\begin{subfigure}[b]{0.32\columnwidth}
	    \includegraphics[width=0.8\textwidth,center]{Figs/Frames2_orientation.pdf}
	    \caption{\scriptsize orientation alignment}
	    \label{fig:frame-b}
	\end{subfigure}
	\begin{subfigure}[b]{0.32\columnwidth}
	    \includegraphics[width=0.8\textwidth,right]{Figs/Frames2_none.pdf}
	    \caption{\scriptsize no alignment}
        \label{fig:frame-c}
	\end{subfigure}
\caption{\edit{Required alignment of UAV frames in existing swarm strategies: (a) the most restrictive case requiring a common reference frame, i.e., orientation and origin of the frames must be aligned; (b) only the orientation of the frames must be aligned; (c) no alignment restrictions (this work).}}
	\label{fig:Frames}
\end{figure}




\subsection{Related Work}

Existing aerial swarms can be grouped based on the coordination (centralized vs.\ distributed) and localization (external vs.\ onboard) methods used. 
\edit{It is further crucial to distinguish these methods based on the level of alignment required for the vehicle coordinate frames; see Fig.~\ref{fig:Frames}.} 
 
\edit{
Works with \textit{centralized} coordination and \textit{external} localization include~\cite{Preiss2017, Honig2018, Du2019}, which are based on lightweight UAVs with limited onboard computational capability and therefore rely on an external motion capture system and a base station.
Works with \textit{distributed} coordination and \textit{external} localization include \cite{wilson2020robotarium}, \cite{enright2004spheres}, where robots execute distributed controls  based on external localization by motion capture and ultrasonic beacons, respectively.
Works with \textit{centralized} coordination and \textit{onboard} localization include~\cite{Forster2013}, \cite{Loianno2016}, which use a ground station for task assignment among vehicles.
In \cite{Weinstein2018}, formation flying based on VIO is demonstrated, where motion planning and assignment are run on a base station to ensure collision-free trajectories.
The coordination strategies used in aforementioned works require a \textit{common reference frame} (Fig.~\ref{fig:frame-a}).
}


\edit{
Despite the large body of work on formation control~\cite{Oh2015}, and the variety of onboard sensing solutions for localization (e.g., VIO~\cite{Delmerico2018}), few frameworks demonstrated formation flying with \textit{distributed} coordination and \textit{onboard} localization.
A key reason is reliance of many distributed control and assignment algorithms on aligned frames (Fig.~\ref{fig:frame-a}, \ref{fig:frame-b}), which require computation-expensive and/or communication-intensive synchronization/consensus steps for frame alignment.
Equally important, dependence on alignment in existing methods \cite{Wang2017,Turpin2014, van2011reciprocal, morgan2016swarm} diminishes robustness to inherent noise and unobservable errors that cannot be corrected (e.g., disparities between the actual and estimated body frame \textit{orientation} caused by VIO drift).
Leveraging coordination methods that are \textit{robust to misaligned frames} is hence crucial and a focus of this work. 
}






\edit{
Examples of other pipelines with distributed coordination and onboard localization include \cite{Montijano2016,Tron2016}.
Both works demonstrated formation flying on three UAVs, required information from an external motion capture system due to hardware limitations, did not incorporate collision avoidance, and required frame alignment.
}
\edittwo{Note that while~\cite{Montijano2016,Tron2016} can achieve formations with arbitrary headings as illustrated in Fig.~\ref{fig:frame-c}, knowledge of relative orientations is still required; therefore, they belong to the category of Fig.~\ref{fig:frame-b}.}






\if 0

\r{
decentralized coordination setting combined with VIO:
D-CAPT [26]~\cite{}:
ORCA ~\cite{}: 
CBF [2]~\cite{} :
[A]
}

\r{Robusteness in coordination,  with compounded noise/latency, which would eventually break (b).\\


some existing algorithm might as well
work in a similar fully decentralized setting, when combined with VIO
as proposed here. For example, D-CAPT [26], ORCA, CBF [2] might also be
useful for such a task and are computationally even more efficient than
the proposed approach. \\

R2:  onboard sensing for localization ->
 Finally, the related work section only
focuses on this aspect of the pipeline, discussing how many formation papers include
onboard localization but barely sells the advantages of the coordination module (the actual
proposal of the paper) against other competitors such as [26] or [A] or to mention similar
coordination pipelines. \\


Given a solution to this problem, the controller in Section III seems unnecessary, each drone
has a target position and can use a local controller with collision avoidance that drives it to
that position. Note that such controllers exists in the literature (e.g., RVO in any of its
multi-agent variantes), they are distributed in nature and only require local sensing.


}

\fi




\begin{figure*}
  \centering
    \includegraphics[width=\textwidth]{images/overview.png}
  \caption{
  Overview of our approach. (a) We aim to find a collision-free path for a rigid body from its current configuration $q_t$ to the goal configuration $q_g$. (b) We assume no prior knowledge about the scene and represent obstacles by points and normals sampled on object surfaces. (c) Our neural network learns the PointNet encoding of observed points and normals together with the motion policy. (d) The learned network generates actions that move the body towards the goal configuration along a collision-free path.   
 }
  \label{fig:overview}
\vspace{-0.4cm}
\end{figure*}



\section{Distributed Formation Control} \label{sec:motionplanning}

\edit{
To achieve a given desired formation, we require a distributed strategy in which the UAVs execute their motions independently and are robust to misalignments of UAV frames.}
To make the paper self-contained, we first review a candidate strategy and then present a solution to address the scalability issue that arises for large-scale formations.


\subsection{Overview of Formation Control}

The framework for achieving a formation using only relative and local position measurements is based on \cite{Lin2016, Lin2016a} and our previous work \cite{Fathian2019, Fathian2018b}. 
For each UAV, the key steps in this strategy can be summarized as follows.
\begin{enumerate}
\item
The UAV calculates the position vectors from itself to each of its neighbors in its own body frame.
\item Each vector is scaled and rotated about the $z$-axis of the UAV's body frame. The amount of scaling and rotation is pre-specified and depends on the desired formation. 
\item These scaled and rotated vectors are then summed to obtain a resultant velocity vector command.
\end{enumerate}

\edit{We emphasize that the above strategy does \textit{not} rely on a common reference frame  (Fig.~\ref{fig:frame-a}) and the scaling and rotation are performed in each respective UAV body frame.} To formulate and analyze this framework mathematically, however, we consider a common reference frame, in which we express the position of UAV $i$ by $q_i \in \br^3$ and the  vector connecting UAV $i$ to its neighbor $j \in \mathcal{N}_i$  by $q_j - q_i$. The scaling and rotation of this vector is expressed as $A_{ij} (q_j - q_i)$, where $A_{ij} \in \ba(3)$ is called a \textit{gain matrix} and belongs to the set of scaled rotation matrices along the $z$-axis denoted by 
\begin{equation} \label{eq:Aset}
\ba(3) \eqdef \left\{ \begin{bsmallmatrix}
a & -b & \; 0 \\
b & ~a & \; 0 \\
0 & ~0 & \; c \\
\end{bsmallmatrix} \,:\, a,b,c \in \br  \right\}.
\end{equation}
Consequently, the motion of UAV $i$ can be expressed as 
\begin{equation} \label{eq:control}
\dot{q}_i = \sum_{j \in \mathcal{N}_i}{A_{ij} \, (q_j - q_i)},
\end{equation}
where $\dot{q}_i$ is the velocity vector that encapsulates the desired speed and direction of motion for the vehicle.
\edittwo{While we consider single-integrator dynamics for simplicity of motion planning, higher-order dynamics can be utilized~\cite{Fathian2018b}.}


In \eqref{eq:control}, we assume that the $z$-axes of UAVs' body frames (and the reference frame used for the analysis) are aligned. In practice, the direction of gravity can be used to align these axes, and, as we will discuss in Section~\ref{sec:robustness}, small misalignments caused by measurement errors or acceleration effects do not affect the convergence. 
\edit{
Note that we do not require that the $x$-$y$ axes be aligned; the UAVs can have arbitrary yaw orientations (Fig.~\ref{fig:frame-c}). This point distinguishes \eqref{eq:control} from the consensus-based \cite{Ren2007}, bearing-based \cite{Zhao2019}, or similar distributed control \cite{Montijano2014}, in which convergence guarantees rely on orientation alignment or consensus of the UAV body frames (Fig.~\ref{fig:frame-b}). }

To analyze the trajectory of the swarm, we define  
\begin{equation} \label{eq:qA}
q \eqdef \begin{bsmallmatrix}
q_{1} \\
q_{2} \\
\vdots \\
q_{n}   
\end{bsmallmatrix}, \quad 
A \eqdef \begin{bsmallmatrix}
-\sum_{j} A_{1j} & A_{12} & \cdots & A_{1n} \\
A_{21} & -\sum_{j} A_{2j} & \cdots & A_{2n} \\
\vdots &                          & \ddots & \vdots \\
A_{n1} &       A_{n2}             & \cdots &  -\sum_{j} A_{nj}   
\end{bsmallmatrix},
\end{equation}
where $q$ is the aggregate vector of UAV positions and $A$ consists of gain matrices. Here, if UAVs $i$ and $j$ are not neighbors, the corresponding $A_{ij}$ is defined as a zero matrix.
Based on \eqref{eq:control} and by using the notation in \eqref{eq:qA}, swarm motion can be expressed by $\dot{q} = A\, q$, which determines the trajectories that the UAVs traverse.


Given a desired formation expressed via the set of points ${p_i = [x_i,\, y_i, \, z_i]^\top} \in \br^3$, we define
\begin{equation} \label{eq:N}
N \eqdef \begin{bsmallmatrix}
p^x_1 & p^y_1 & p^z_1 &  e^x & e^y & e^z \\ 
p^x_2 & p^y_2 & p^z_2 &  e^x & e^y & e^z \\
\vdots &  \vdots & \vdots & \vdots & \vdots & \vdots \\ 
p^x_n & p^y_n & p^z_n &  e^x & e^y & e_z 
\end{bsmallmatrix} \in \mathbb{R}^{3n\times 6},
\end{equation}
where ${p^x_i \eqdef [x_i,\, y_i,\, 0]^\top}$, ${p^y_i \eqdef [-y_i,\, x_i,\, 0]^\top}$, ${p^z_i \eqdef [0,\, 0,\, z_i]^\top}$, ${e^x \eqdef [1,\,0,\,0]^\top}$, ${e^y \eqdef [0,\,1,\,0]^\top}$, and ${e^z \eqdef [0,\,0,\,1]^\top}$. 
Convergence to the desired shape is guaranteed if $A\, N = 0$ (i.e., columns of $N$ are null vectors of $A$) and if all remaining eigenvalues of $A$ not associated with $N$ are strictly negative.
For such an $A$ to exist, each UAV should have a sufficient number of neighbors. Specifically, if vertices and edges represent the UAVs and their neighboring relations in the formation graph $\sg$, $A$ exists if $\sg$ is universally rigid (see~\cite[Thm. 3.2]{Lin2016a}).

\edit{Our approach consists of \textit{design} and \textit{execution} phases.
In the \textit{design} phase, a gain matrix $A$ as in~\eqref{eq:qA} is computed offline based on the specification of a formation \edittwo{graph $\sg$} and serves as an input to the distributed algorithms that run onboard the UAVs.
The \textit{execution} phase is entirely distributed, where the UAVs plan their trajectories independently using relative translation measurements to their neighboring UAVs.}







\subsection{Scalable Gain Design}

Given a desired formation, the gain matrix $A$ that meets the aforementioned constraints can be computed from
\begin{equation} \label{eq:OptimCVX}
\begin{aligned}
& \underset{A \in \bs^{-}_{3n}}{\text{minimize}}
& & \lambda_{\text{max}} \left( Q^\top A \, Q \right) & \\
& \text{subject to}
& & A\, N = 0 & \\
&&& A_{ij} \in \mathbb{A}(3)  & \forall_{i,j} \\
&&& A_{ij} = 0 \;   & \forall_{i} \ \forall_{j \notin \mathcal{N}_i} \\
&&& \tr(A) = \text{constant}
\end{aligned}
\end{equation}
where $\lambda_{\text{max}}$ denotes the largest eigenvalue of a matrix, $Q \in \mathbb{R}^{3n\times (3n-6)}$ is the orthogonal complement of $N$ (i.e., $N^\top Q = 0$), which is found from the singular value decomposition of $N$, and $\bs^-_{3n}$ is the space of symmetric negative semidefinite matrices of dimension $3n$.
The objective of \eqref{eq:OptimCVX} is to make the nonzero eigenvalues of $A$ as negative as possible ($Q^\top A \, Q$ is the restriction of $A$ on the subspace $Q$ and eliminates the zero eigenvalues associated with $N$).
By doing so, stability and robustness of the formation to noise, measurement errors, and disturbances is increased.  
We note that the last constraint in \eqref{eq:OptimCVX} sets the trace of $A$ to a constant value to ensure that the problem is bounded (without this constraint, if $A\in\bs^-_{3n}$ is a solution, so is $c\, A$ for any $c>0$ with a better objective value).
The universal rigidity assumption on the formation graph \cite[Thm. 3.2]{Lin2016a} is sufficient to ensure that \eqref{eq:OptimCVX} is feasible and that all remaining eigenvalues of $A$ not associated with $N$ are strictly negative.


The formulation \eqref{eq:OptimCVX} was presented in our earlier work~\cite{Fathian2019} and can be solved relatively quickly using existing SDP solvers for small number of vehicles. However, for large-scale problems (e.g., more than $50$ UAVs) it becomes challenging, or even impossible, to solve. 
We address this issue by exploiting the problem structure to derive a solution based on the alternating direction method of multipliers (ADMM).



We observe from \eqref{eq:Aset} that $A_{ij}\in\ba(3)$ has a block diagonal structure, which can be expressed by ${A_{ij} = \mathrm{blkdiag}(D_{ij}, c_{ij})}$, where the $2\times 2$ matrix $D_{ij}$ consists of the first two rows and columns, and the scalar $c_{ij}$ is the entry in the last row and column of $A_{ij}$.
Due to this structure, we conclude from \eqref{eq:control} that vehicle trajectories along the $x$-$y$  and $z$ components are decoupled and depend only on $D_{ij}$ and $c_{ij}$, respectively.
This observation allows us to split \eqref{eq:OptimCVX} into two subproblems with lower dimensions, leading to reduced computational effort.
By defining
\begin{equation} \label{eq:} 
B \eqdef \begin{bsmallmatrix}
-\sum_{j} c_{1j} & c_{12} & \cdots & c_{1n} \\
c_{21} & -\sum_{j} c_{2j} & \cdots & c_{2n} \\
\vdots &                      & \ddots & \vdots \\
c_{n1} &     c_{n2}       & \cdots &  -\sum_{j} c_{nj}   
\end{bsmallmatrix}, ~~
M \eqdef \begin{bsmallmatrix}
z_1 & 1 \\ 
z_2 & 1 \\
\vdots &  \vdots \\ 
z_n & 1
\end{bsmallmatrix},
\end{equation}
which correspond to the $z$ components of $A$ in \eqref{eq:qA} and $N$ in \eqref{eq:N}, the problem of finding $c_{ij}$ is formulated as
\begin{equation} \label{eq:OptimZ}
\begin{aligned}
& \underset{B \in \bs^{-}_{n} }{\text{minimize}}
& & \lambda_{\text{max}} \left(  R^\top B \, R \right) & \\
& \text{subject to}
& & B \, M = 0 & \\
&&& c_{ij} = 0 \;   & \forall_{i} \ \forall_{j \notin \mathcal{N}_i} \\
&&& \tr(B) = \text{constant}
\end{aligned}
\end{equation}
where $R \in \br^{n \times (n-2)}$ is the orthogonal complement of $M \in \br^{n\times 2}$. 
The optimization problem for finding $D_{ij}$ is formulated similarly to \eqref{eq:OptimZ}, with an additional constraint that the diagonal entries of $D_{ij}$ must be equal and the off-diagonal entries must have the same absolute value with different signs. With this point in mind, we henceforth focus our attention on \eqref{eq:OptimZ}. 
The following proposition brings \eqref{eq:OptimZ} into the standard form suitable for applying ADMM.


\begin{prop} \label{prop:SDP}
Problem \eqref{eq:OptimZ} can be formulated as
\begin{equation} \label{eq:OptimX}
\begin{aligned}
& \underset{X \in \bs^+_{2n-4} }{\text{minimize}}
& & \left< C,\, X \right> & \\
& \text{subject to}
& & \sa(X) = b  
\end{aligned}
\end{equation}
where $\left< C,\, X \right> \eqdef \tr(C^\top X)$,  
$C \eqdef \begin{bsmallmatrix}
I & 0 \\
0 & 0
\end{bsmallmatrix}$ with $I$ as the identity matrix of size $n-2$, and $b \in \br^m$. The operator $\sa(X)$ represents a set of linear constraints on $X$ and enforces it to have the block diagonal structure $X \eqdef \begin{bsmallmatrix}
\gamma\, I && I \\
I && Z 
\end{bsmallmatrix}$, where $\gamma \geq 0$ and $Z \in \bs^+_{n-2}$. The solution of \eqref{eq:OptimZ} is obtained from $X$ as $B = - M \, Z \, M^\top$.

\end{prop}

\edit{Proposition~\ref{prop:SDP} is proved in the appendix of~\cite{lusk2020distributed}.} 
We now leverage the ADMM technique in \cite{Wen2010} to solve  \eqref{eq:OptimX}.
From~\cite{Wen2010}, the augmented Lagrangian associated with \eqref{eq:OptimX} is 
\begin{equation} \label{eq:Lagrangian}
\mathcal{L} \eqdef - \langle y, b \rangle + \langle \sa^*(y) + S - C,\, X \rangle + \frac{1}{2\mu} \| \sa^*(y) + S - C \|,
\end{equation}
where $S \in \bs^+_{2n-4}$ and $y$ are dual variables associated with constraints $X \in \bs^+_{2n-4}$ and $\sa(X) = b$, respectively, $\sa^*$ is the adjoint of $\sa$, and $\mu > 0$ is a penalty parameter that balances the standard Lagrangian and the augmented term. 
ADMM then proceeds by alternatively optimizing each primal and dual variable with others fixed, which results in a closed-form solution for each subproblem. Denoting by the superscript $k$ the iteration number, the ADMM iterative update procedure is given as
\begin{equation} \label{eq:ADMM}
\begin{aligned}
y^{k+1} &= (\sa \sa^*)^{-1} \left( \sa (C - S^k - \mu X^k) + \mu\, b \right), \\
W^{k+1} &=  C - \sa^* (y^{k+1}) - \mu \, X^k, \\
S^{k+1} &= \mathcal{P}_{\mathrm{psd}} \left(W^{k+1} \right), \\
X^{k+1} 
&= \frac{1}{\mu} \left( S^{k+1} - W^{k+1} \right).
\end{aligned}
\end{equation}
In~\eqref{eq:ADMM}, the operator $\mathcal{P}_{\mathrm{psd}}$ denotes the projection onto the positive semidefinite cone $\bs^+$, and is computed via  eigendecomposition (see \cite{Wen2010} for details).
ADMM typically converges in a reasonable time to a solution with acceptable accuracy.
The number of ADMM iterations required for convergence depends on the desired accuracy as well as the formation graph (e.g., when the formation graph is complete, it is straightforward to show that ADMM converges to the optimal solution in a single iteration).
The time comparisons between an existing SDP solver for \eqref{eq:OptimCVX} and the presented ADMM method \eqref{eq:ADMM} are given in Section~\ref{sec:exper}. 

\subsection{Robustness, Collision Avoidance, and Formation Size}\label{sec:robustness}



Gain matrices are recomputed only when a new desired formation is specified.
During execution, vehicles use the gains and the relative position of their neighbors to compute the velocity vector $u_i$ in \eqref{eq:control} at each time instance. Having $u_i$ computed, the vehicle's low-level controller is tasked with tracking the direction and speed specified by this vector.

  



One can show that 1) \textit{any} \textit{positive} scaling; and 2) \textit{any} rotation less than \textit{90 degrees} of the velocity vector $u_i$ does not void the convergence guarantees of the formation control strategy (see \cite[Thm. 2]{Fathian2019}). 
These key properties indicate extreme robustness to errors and disturbances. For example, discrepancies between the actual and desired velocity of a vehicle caused by imperfect tracking, unmodeled dynamics, or small misalignments in $z$-axes of UAV body frames can be modeled as a positive scaling and small rotation of the nominal $u_i$, for which  convergence to the desired formation is unaffected.
The aforementioned properties can be further used for collision avoidance, where velocity vectors that lead vehicles to close proximity are modified to prevent collisions. More specifically, \eqref{eq:control} can be modified as
\begin{gather} \label{eq:RotationCtrl}
u_i \eqdef c_i \, R_i \sum_{j \in \mathcal{N}_i}  A_{ij} \, (q_j - q_i),
\end{gather}
where the rotation matrix $R_i$, which is limited to $90$ degrees, is chosen to rotate any velocity vector that brings two vehicles closer than a specified distance. 
If there is no feasible direction of motion within this range, the scalar $c_i$, which is normally set to one, is set to zero to stop the vehicle.



The collision avoidance strategy \eqref{eq:RotationCtrl} runs onboard, but comes at the cost of losing convergence guarantees since vehicles can become gridlocked due to the unavailability of motion directions (allowing $c_i = 0$ in \eqref{eq:RotationCtrl} violates the aforementioned property in which $c_i > 0$ is required to ensure convergence). 
Optimal assignment of vehicles to target formation points guarantees non-intersecting lines from their current positions to their assigned points (see \cite[Thm. 3.1]{Turpin2014}).
\edit{Since the vehicle body frames can be nonaligned and the control is distributed, vehicles are not expected to move perfectly in a straight line under our control strategy. However, assignment can still help deconflict the swarm by increasing the availability of motion directions.
The impact of assignment on resolving gridlocks is shown in Section~\ref{sec:exper}. }





Finally, note that \eqref{eq:RotationCtrl} leads to achieving the desired formation \textit{shape}, but the formation \textit{size} is not regulated and depends on the initial position of the vehicles. To control the size, \eqref{eq:RotationCtrl} can be augmented to contract (or expand) the formation when the vehicles are farther (or closer) than the desired distance. 
This augmented strategy, and its theoretical convergence guarantees, is discussed in our earlier work for 2D formations (see (74) in \cite{Fathian2018b}). Since the extension to 3D formations considered in this work is straightforward, we omit this discussion for brevity. 













\section{\edit{Distributed} Task Assignment} \label{sec:assign}


The goal of task assignment is to uniquely allocate each UAV to a point in the desired formation.
A natural objective for this task is to minimize the overall distance from the UAVs to their assignments in the desired formation.
We are interested in the final 3D geometric shape rather than the exact location and yaw of the end formation.
To this effect, we allow a rotation $R$ around the $z$-axis, and translation $t$ of the desired formation coordinates that minimize the overall distance from the UAVs to the rotated and translated desired formation.
More precisely, our objective is to find the assignment map $\sigma$ that solves
\begin{equation}\label{eq:assign_prob}
\begin{aligned}
\underset{\substack{R\in\mathcal{R}_z,\;t\in\br^3\\\sigma\in\mathrm{S}_n}}{\text{minimize}} && \sum_{i=1}^{n} \| q_i - (R \, p_{\sigma(i)}+t) \|^2,
\end{aligned}
\end{equation}
where $q_i$ denotes the UAV positions, 
$\mathrm{S}_n$ is the symmetric group of all permutations from the set $\{1,\dots, n\}$ to itself,
and $\mathcal{R}_z$ is the set of rotation matrices around the $z$-axis. Recall that the 
$z$-axes of UAV body frames are assumed to be aligned, as per Section~\ref{sec:motionplanning}.


\edit{Three challenges arise in finding a distributed solution for~\eqref{eq:assign_prob}.
First, the objective of \eqref{eq:assign_prob} includes the positions of all UAVs, whereas each UAV only obtains the positions of its neighbors.
Second, the UAV body frames are not aligned and the UAVs only know the transformations between their body frames and their neighbors.  
Lastly, \eqref{eq:assign_prob} is a non-convex mixed integer program, for which finding the global optimizer becomes intractable for large $n$.} 
As computational efficiency and scalability are of utmost concern for UAV platforms, 
we settle with obtaining a suboptimal answer via a coordinate descent approach and an auction assignment strategy.
This approach is inspired by \cite{macdonald2011multi}, which in contrast uses a centralized Hungarian algorithm for assignment.
Every iteration of our algorithm consists of an \textit{alignment} stage and an \textit{assignment} stage, where the assignment is fixed as we solve for an alignment, and vice versa.


\subsection{Alignment}

Given an assignment $\sigma^*$ (e.g., identity assignment ${\sigma^*(i)=i}$ for every new formation, or prior assignment computed for the same formation), UAV $i$ solves a \edit{distributed} formulation of \eqref{eq:assign_prob} given by
\begin{equation}\label{eq:assign_prob_fix_sigma}
\begin{aligned}
\underset{\substack{R_i\in \mathcal{R}_z,\;t_i\in\br^3}}{\text{minimize}} && \sum_{j\in\mathcal{N}'_i} \| q_j - (R_i \, p_{\sigma^*(j)} + t_i) \|^2,
\end{aligned}
\end{equation}
where ${\mathcal{N}}'_i \eqdef \mathcal{N}_i\cup \{i\}$, and positions $q_j$ are in UAV $i$'s start frame.
In \eqref{eq:assign_prob_fix_sigma}, UAV $i$ aims to align the desired formation to minimize the distance to its own and neighbors' positions based on the given assignment $\sigma^*$.   
Fig.~\ref{fig:alignment} gives an illustrative example of this stage.
Problem \eqref{eq:assign_prob_fix_sigma} is the well-known point cloud alignment problem, for which the optimal solution $\left(R^*_i,t^*_i\right)$ 
is obtained from Arun's method~\cite{arun1987least}  using the projection of $q_j$ and $p_{\sigma^*(j)}$ on the {$x$-$y$ plane} for the rotation. %



\begin{figure} [t!]
	\centering
	\includegraphics[width=\columnwidth]{Figs/alignment4.pdf}    
		\caption{
		Illustrative 2D alignment example with four vehicles from UAV~1's perspective.
        UAV $4$ is not shown because it does not communicate with UAV $1$.
		\textbf{(left)}~New formation graph, UAV $1$ and its neighbors before the alignment.
		\textbf{(right)}~Aligned formation 
		based on UAV $1$, its neighbors and their corresponding formation points. 
		The formation point associated to UAV 4 is faded to indicate that UAV 1 does not have information about it.
		}	
	\label{fig:alignment}
\end{figure}






\subsection{Assignment}

In this stage, the UAVs aim to collaboratively find an assignment based on the results obtained from their alignment stage. The assignment problem is formulated as
\begin{equation}\label{eq:assign_prob_fix_Rt}
\begin{aligned}
\underset{\substack{\sigma\in\mathrm{S}_n}}{\text{minimize}} && \sum_{i=1}^{n} \| q_i - (R_i^* \, p_{\sigma(i)}+t_i^*) \|^2.
\end{aligned}
\end{equation}
Problem~\eqref{eq:assign_prob_fix_Rt} is a linear sum assignment problem that can be solved optimally by methods such as the distributed Hungarian algorithm~\cite{giordani2010distributed, chopra2017distributed}, which converges in $\text{O}(n^3)$ iterations. 
Due to onboard resource constraints, we trade optimality for computational efficiency by using our prior work, the consensus-based auction algorithm (CBAA)~\cite{choi2009consensus}. 
CBAA is guaranteed to converge in at most $n d$ iterations, where $d$ is the diameter of the formation graph, $\mathcal{G}$.




To bring \eqref{eq:assign_prob_fix_Rt} into the standard form for applying CBAA, let binary variables $x_{ij}$ represent the assignment $\sigma$, where ${x_{ij} = 1}$ if ${\sigma(i) = j}$ and $0$ otherwise.
Further, let ${c_{ij} \eqdef 1\,/\,\| q_i-(R_i^* \, p_j + t_i^*) \|^2}$ denote the positive score for assigning UAV $i$ to formation point $j$. In practice, a small positive number can be added to the denominator of $c_{ij}$ to avoid division by zero.
It is straightforward to show that~\eqref{eq:assign_prob_fix_Rt} can be expressed as the integer program
\begin{equation}\label{eq:cbaa_problem}
\begin{aligned}
\underset{x_{ij} \in \{0,1\}}{\text{maximize}} &
&&\sum_{i,j=1}^{n} c_{ij} \,x_{ij}, \\
\text{subject to}
&&& \textstyle\sum_{i=1}^{n} x_{ij} = 1, \; \forall_{j} \\
&&& \textstyle\sum_{j=1}^{n} x_{ij} = 1,  \; \forall_{i} 
\end{aligned}
\end{equation}
where the constraints on $x_{ij}$ enforce conflict-free and one-to-one assignment captured by $\sigma \in S_n$ in~\eqref{eq:assign_prob_fix_Rt}, and maximizing~\eqref{eq:cbaa_problem} is equivalent to minimizing the overall distance from the UAVs to the rotated and translated formation in~\eqref{eq:assign_prob_fix_Rt}.





    



In executing CBAA, UAV $i$ stores and updates its own assignment and a list of winning bids (initialized as zeros) for all formation points. Each iteration of CBAA consists of an auction phase and a consensus phase. 
In the auction phase, UAV $i$ determines which formation point it would like to be assigned to in three steps: (1) check if any formation point $p_j$ produces a score $c_{ij}$ higher than its current winning bid; (2) of those formation points, set $x_{ij}=1$ for the $p_j$ that produces the highest score; (3) update the bid for the winning $p_{j}$ with new score $c_{ij}$.
In the consensus phase, vehicles converge on a common winning bid list. UAV $i$ exchanges its winning bid list with its neighbors and updates its list with the highest values from its own and all received lists. It sets $x_{ij}=0$ if the new winning bid for $p_j$ is higher than $c_{ij}$, implying that a different vehicle has been assigned to $p_j$.

Since CBAA is distributed, no central authority exists to affirm convergence.
Therefore, we enforce a synchronous execution to terminate the algorithm in $n d$ iterations, which is the maximum number of iterations required to guarantee convergence. 
The final assignment is recovered by letting $\sigma^*(i)=j$ for each $x_{ij}=1$.
CBAA guarantees a conflict-free assignment even though UAVs \edittwo{do not have a common reference frame} and may have inconsistent position estimates or different $R^*_i$ and $t^*_i$ for alignment.
\edittwo{We emphasize that although $c_{ij}$ is calculated by each UAV using only local knowledge, CBAA assigns UAVs to formation points without conflict.}
Further, it retains at least $50\%$ of the optimal performance; that is, given the optimal overall score $C^*$ of \eqref{eq:cbaa_problem} and the $C$ resulted from CBAA, $C/C^* \geq 0.5$. 



%!TEX root = main.tex

\section{Experimental Evaluation}
\seclabel{experiments}

We first evaluated our algorithms
in an offline setting~(\secref{offline-expr}), where we record execution traces and evaluate different approaches on the \emph{same} input.
This eliminates biases due to non-deterministic thread scheduling.
Next, we consider an online setting~(\secref{online-expr}),
where we instrument programs and perform the analyses during runtime.
We conducted all our experiments on a standard laptop with \SI{1.8}{GHz} Intel Core i7 processor and \SI{16}{GB} RAM.

% We evaluated our algorithms in two experimental settings. 
% The first setting is offline experiments~(\secref{offline-expr}),
% in which we record execution traces and evaluate different approaches.
% This has the benefit that different approaches can be compared on the \emph{same} input,
% thereby eliminating biases due to non-deterministic thread scheduling.
% The second setting is online experiments~(\secref{online-expr}),
% in which we instrument programs and perform the analyses during runtime.
% We conducted all our experiments on a standard laptop with 1.8GHz Intel Core i7 processor and 16GB RAM.

%We compare our predictive algorithm with the \dlfuzzer tool. 
%This setting enables us to assess the applicability of our technique in a runtime monitoring system.
%The state-of-the-art deadlock predictors \dirk and \seqc work offline by design.
%Hence, they are not applicable for a comparison in this setting.
%As discussed in~\secref{otf}, offline methods are not directly applicable in runtime monitoring systems.

%Talk about implementation - name of tool + prog lang + traces are logged + filtering phase + cycle detection phase + online vs offline + trace conversion for seqc
%
%Setup - benchmarks + log traces using so-and-so-tools + 1-trace-per-benchmark  + cluster details + how many runs per trace + Timeout (if any)
%
%benchmark details: number of shared locks
%\subsection{Evaluation}
%
%Some suggested experiments:
%
%\begin{itemize}
%	\item Comparison with other tool(s) - \dirk and \seqc
%	\item Comparison of online vs online algorithms
%	\item Scalability with number of events
%	\item Scalability with number of threads in the trace
%	\item Scalability with number of threads in the deadlock pattern
%	\item Time spent in each kinds of events (this will be useful to guide future research)
%\end{itemize}
%
%\Andreas{For matching reports, we slightly adapt the algorithm to report all sets of program locations that contain a deadlock pattern (hence we might have more than one reports per abstract deadlock pattern, corresponding to different program locations)}


% Here we report on an implementation and experimental evaluation of our algorithms.

% \subsection{Experimental Setup}



\subsection{Offline Experiments}
\seclabel{offline-expr}


\Paragraph{Experimental setup}
The goal of the first set of experiments is to evaluate 
$\SyncPDOffline$, and compare it
against prior algorithms for dynamic deadlock prediction.
In order for our evaluation to be precise we evaluate all algorithms on the \emph{same} execution trace.
We implemented $\SyncPDOffline$ in Java inside the \toolname analysis tool~\cite{rapid}, 
following closely the pseudocode in \algoref{offline}.
\toolname takes as input execution traces, as defined in \secref{prelim}.
These also include fork, join, and lock-request events.
We compare $\SyncPDOffline$ with two state-of-the-art, 
theoretically-sound albeit computationally more expensive, deadlock predictors,
\seqc~\cite{Cai2021} and \dirk~\cite{Kalhauge2018}, both of which also work on execution traces.




On the theoretical side, the complexity of \seqc is $\Otilde(\NumEvents^4)$, 
as opposed to the $\Otilde(\NumEvents)$ complexity of $\SyncPDOffline$. 
Moreover, \seqc only predicts deadlocks of size $2$, and though it could be extended to handle deadlocks of any size, this would degrade performance further.
\seqc may miss sync-preserving deadlocks even of size $2$, 
but can  detect deadlocks that are not sync-preserving.
Thus \seqc and $\SyncPDOffline$ are theoretically incomparable in their detection capability.
We refer to\begin{pldi}~\cite{arxiv}\end{pldi}\begin{arxiv}~\appref{incomp}\end{arxiv} for examples.
We noticed that \seqc fails on traces with non-well-nested locks --- we encountered one such case in our dataset.
\dirk's algorithm is theoretically complete, i.e., it can find all predictable deadlocks in a trace.
In addition, it can find deadlocks beyond the predictable ones, by reasoning about event values.
However, \dirk relies on heavyweight SMT-solving and
employs windowing techniques to scale to large traces. 
Due to windowing, it can miss deadlocks between events that are outside the given window. 
%\hunkar{
As with previous works~\cite{Cai2021, Kalhauge2018}, we set a window size of $10$K for \dirk.
%}

Our dataset consists of several benchmarks 
from standard benchmark suites --- IBM Contest suite~\cite{Farchi03}, Java Grande suite~\cite{Smith01},
DaCapo~\cite{Blackburn06}, and
SIR~\cite{doESE05} ---
and recent literature~\cite{Kalhauge2018, Cai2021, jula2008deadlock, Joshi2009}.
Each benchmark was instrumented with RV-Predict~\cite{rvpredict} or Wiretap~\cite{Kalhauge2018} and
executed in order to log a single execution trace.

%!TEX root = ../main.tex

\begin{table}[h!]
\caption{
%\hcomment{Should we remove some not very important rows from this table?}
Trace characteristics, abstract lock graph statistics and performance comparison.
%Column 1 denotes the name of the benchmark.
Columns 2-6 show the number of events, threads, variables, locks
and total number of lock acquire and request events.
Columns 7-9 show the number of cycles, abstract and concrete deadlock patterns in the abstract lock graph.
Columns 10 - 15 show the number of deadlocks reported and the times (in seconds) taken. 
by \dirk, \seqc, and \SyncPDOffline.
%Statistics on the lock graph $\lkevgraph{\tr}$ of each trace $\tr$.
% We denote by $\mathcal{N}$, $\mathcal{T}$, $\mathcal{M}$, $\mathcal{L}$ and $\NumAcquires + \mathcal{R}$ the total number of events, number of threads, number of memory locations, number of locks and number of acquire and request events in the benchmarks, respectively. 
%All times are in seconds. 
Time out (T.O) was set to $3$h.
F stands for technical failure.
\label{tab:expr-results}
}
\vspace{-0.17cm}
\setlength\tabcolsep{3pt}
\renewcommand{\arraystretch}{0.91}
\centering
\scalebox{0.86}{
\begin{tabular}{|r|c|c|c|c|c||c|c|c||c|c||c|c||c|c|}
\hline
1 & 2 & 3 & 4 & 5 & 6 & 7 & 8 & 9 & 10 & 11 & 12 & 13 & 14 & 15 \\
\hline
\multirow{2}{*}{\textbf{Benchmark}}& 
\multirow{2}{*}{$\mathcal{N}$} & \multirow{2}{*}{$\mathcal{T}$} & \multirow{2}{*}{$\mathcal{V}$} & \multirow{2}{*}{$\mathcal{L}$} & \multirow{2}{*}{$\NumAcquires/\mathcal{R}$} 
& \multicolumn{3}{c||}{ \textsf{A. Lock Graph }}
& \multicolumn{2}{c||}{ {\dirk} } 
& \multicolumn{2}{c||}{ {\seqc}} 
& \multicolumn{2}{c|}{ {\textsf{$\SyncPDOffline$}}} \\
\cline{7-15}
& & & & & 
& \textsf{|$\textsf{Cyc}$|}
& \textsf{\textsf{A. P.}}
& \textsf{\textsf{C. P.}}
& \textbf{Dlk} 
& {\textbf{Time}} 
&{\textbf{Dlk}} 
& {\textbf{Time}}
&{\textbf{Dlk}} 
&{\textbf{Time}} \\
\hline
Deadlock & 39 & 3 & 4 & 3 & 8 & 1 & 1 & 1 & 1 & 0.02 & 0 & 0.09 & 0 & 0.16\\
NotADeadlock & 60 & 3 & 4 & 5 & 16 & 1 & 1 & 1 & 0 & 0.02 & 0 & 0.09 & 0 & 0.16\\
Picklock & 66 & 3 & 6 & 6 & 20 & 2 & 2 & 2 & 1 & 0.02 & 1 & 0.10 & 1 & 0.18\\
Bensalem & 68 & 4 & 5 & 5 & 22 & 2 & 2 & 2 & 1 & 0.02 & 1 & 0.12 & 1 & 0.16\\
Transfer & 72 & 3 & 11 & 4 & 12 & 1 & 1 & 1 & 1 & 0.02 & 0 & 0.09 & 0 & 0.15\\
Test-Dimmunix & 73 & 3 & 9 & 7 & 26 & 2 & 2 & 2 & 2 & 0.02 & 2 & 0.10 & 2 & 0.17\\
StringBuffer & 74 & 3 & 14 & 4 & 16 & 1 & 3 & 6 & 2 & 0.02 & 2 & 0.12 & 2 & 0.19\\
Test-Calfuzzer & 168 & 5 & 16 & 6 & 48 & 2 & 1 & 1 & 1 & 0.02 & 1 & 0.12 & 1 & 0.17\\
DiningPhil & 277 & 6 & 21 & 6 & 100 & 1 & 1 & 3K & 1 & 1.60 & 0 & 0.09 & 1 & 0.17\\
HashTable & 318 & 3 & 5 & 3 & 174 & 1 & 2 & 43 & 2 & 0.19 & 2 & 0.12 & 2 & 0.19\\
Account & 706 & 6 & 47 & 7 & 134 & 3 & 1 & 12 & 0 & 0.19 & 0 & 0.09 & 0 & 0.18\\
Log4j2 & 1K & 4 & 334 & 11 & 43 & 1 & 1 & 1 & 1 & 0.65 & 1 & 0.11 & 1 & 0.20\\
Dbcp1 & 2K & 3 & 768 & 5 & 56 & 2 & 2 & 3 & - & F & 2 & 0.11 & 2 & 0.19\\
Dbcp2 & 2K & 3 & 592 & 10 & 76 & 1 & 2 & 4 & - & F & 0 & 0.10 & 0 & 0.18\\
Derby2 & 3K & 3 & 1K & 4 & 16 & 1 & 1 & 1 & 1 & 0.23 & 1 & 0.10 & 1 & 0.17\\
RayTracer & 31K & 5 & 5K & 15 & 976 & 0 & 0 & 0 & - & F & 0 & 0.15 & 0 & 0.19\\
jigsaw & 143K & 21 & 8K & 2K & 67K & 172 & 12 & 70 & - & F & 2 & 0.36 & 1 & 1.55\\
elevator & 246K & 5 & 727 & 52 & 48K & 0 & 0 & 0 & 0 & 1.65 & 0 & 0.33 & 0 & 0.27\\
hedc & 410K & 7 & 109K & 8 & 32 & 0 & 0 & 0 & 0 & 2.09 & 0 & 0.50 & 0 & 0.24\\
JDBCMySQL-1 & 442K & 3 & 73K & 11 & 13K & 2 & 4 & 6 & 2 & 28.45 & 2 & 0.24 & 2 & 0.48\\
JDBCMySQL-2 & 442K & 3 & 73K & 11 & 13K & 4 & 4 & 9 & 1 & 3.37 & 1 & 0.22 & 1 & 0.33\\
JDBCMySQL-3 & 443K & 3 & 73K & 13 & 13K & 5 & 8 & 16 & 1 & 31.23 & 1 & 0.25 & 1 & 0.45\\
JDBCMySQL-4 & 443K & 3 & 73K & 14 & 13K & 5 & 10 & 18 & 2 & 5.51 & 2 & 0.28 & 2 & 0.49\\
cache4j & 775K & 2 & 46K & 20 & 35K & 0 & 0 & 0 & 0 & 5.86 & 0 & 0.46 & 0 & 0.39\\
ArrayList & 3M & 801 & 121K & 802 & 176K & 9 & 3 & 672 & 3 & 8.7K & 3 & 21.98 & 3 & 1.68\\
IdentityHashMap & 3M & 801 & 496K & 802 & 162K & 1 & 3 & 4 & 1 & 443.93 & 1 & 8.51 & 1 & 1.45\\
Stack & 3M & 801 & 118K & 2K & 405K & 9 & 3 & 481 & 1 & T.O & 3 & 25.34 & 3 & 2.94\\
Sor & 3M & 301 & 2K & 3 & 719K & 0 & 0 & 0 & 0 & 15.89 & 0 & 44.12 & 0 & 0.61\\
LinkedList & 3M & 801 & 290K & 802 & 176K & 9 & 3 & 10K & 3 & 4.7K & 3 & 48.02 & 3 & 2.06\\
HashMap & 3M & 801 & 555K & 802 & 169K & 1 & 3 & 10K & 3 & 4.4K & 2 & 504.36 & 2 & 1.65\\
WeakHashMap & 3M & 801 & 540K & 802 & 169K & 1 & 3 & 10K & - & T.O & 2 & 499.68 & 2 & 1.70\\
Swing & 4M & 8 & 31K & 739 & 2M & 0 & 0 & 0 & - & F & 0 & 0.72 & 0 & 0.88\\
Vector & 4M & 3 & 15 & 4 & 800K & 1 & 1 & 1B & - & T.O & 1 & 1.52 & 1 & 1.90\\
LinkedHashMap & 4M & 801 & 617K & 802 & 169K & 1 & 3 & 10K & 2 & 40.74 & 2 & 492.87 & 2 & 1.69\\
montecarlo & 8M & 3 & 850K & 3 & 26 & 0 & 0 & 0 & 0 & 2.6K & 0 & 1.81 & 0 & 0.79\\
TreeMap & 9M & 801 & 493K & 802 & 169K & 1 & 3 & 10K & 2 & 105.45 & 2 & 480.11 & 2 & 1.92\\
hsqldb & 20M & 46 & 945K & 403 & 419K & 0 & 0 & 0 & - & F & - & - & 0 & 2.38\\
sunflow & 21M & 16 & 2M & 12 & 1K & 0 & 0 & 0 & - & F & 0 & 8.35 & 0 & 1.62\\
jspider & 22M & 11 & 5M & 15 & 10K & 0 & 0 & 0 & - & F & 0 & 8.49 & 0 & 1.95\\
tradesoap & 42M & 236 & 3M & 6K & 245K & 2 & 1 & 4 & - & F & 0 & 108.16 & 0 & 7.06\\
tradebeans & 42M & 236 & 3M & 6K & 245K & 2 & 1 & 4 & - & F & 0 & 116.23 & 0 & 7.26\\
eclipse & 64M & 15 & 10M & 5K & 377K & 9 & 5 & 280 & - & F & 0 & 26.67 & 0 & 9.90\\
TestPerf & 80M & 50 & 599 & 9 & 197K & 0 & 0 & 0 & 0 & 795.04 & 0 & 47.56 & 0 & 4.30\\
Groovy2 & 120M & 13 & 13M & 10K & 69K & 0 & 0 & 0 & 0 & 1.7K & 0 & 38.06 & 0 & 8.92\\
Tsp & 200M & 6 & 24K & 3 & 882 & 0 & 0 & 0 & 0 & 7.6K & 0 & 72.62 & 0 & 12.70\\
lusearch & 203M & 7 & 3M & 98 & 273K & 0 & 0 & 0 & 0 & 1.3K & 0 & 75.88 & 0 & 14.44\\
biojava & 221M & 6 & 121K & 79 & 16K & 0 & 0 & 0 & - & F & 0 & 63.79 & 0 & 12.65\\
graphchi & 241M & 20 & 25M & 61 & 1K & 0 & 0 & 0 & - & F & 0 & 102.05 & 0 & 25.25\\
\hline\hline\textbf{Totals} & \textbf{1B} & \textbf{7K} & \textbf{70M} & \textbf{37K} & \textbf{8M} & \textbf{256} & \textbf{93} & \textbf{1B} & \textbf{35} & \textbf{>18h} & \textbf{40} & \textbf{2801} & \textbf{40} & \textbf{135}\\\hline
\end{tabular}
}
\end{table}

\Paragraph{Evaluation}
\cref{tab:expr-results} presents our results.
A bug identifies a unique tuple of source
code locations corresponding to events participating in the deadlock.
%Columns 2-6 present the characteristics of our benchmark traces.
Trace lengths vary vastly from $39$ to about $241$M, while the number of threads ranges from $3$ to about $800$,
which are representative features of real-world settings.
\texttt{Hsqldb} contains critical sections that are not well nested, 
and \seqc was not able to handle this benchmark;
our algorithm does not have such a restriction.
%We were unable to run \dirk on certain benchmarks due to \dirk's technical issues, which are marked with F.


\vspace{-0.1cm}
\SubParagraph{\underline{Abstract vs Concrete Patterns}}
Columns 7-9 present statistics on the 
abstract lock graph $\lkevgraph{\tr}$ of each trace $\tr$.
Many traces have a large number of concrete deadlock patterns 
but much fewer abstract deadlock patterns;
a single abstract deadlock pattern can 
comprise up to an order of $10^4$ more concrete patterns (Column $8$ v/s Column $9$).
Unlike all prior sound techniques, 
our algorithms analyze 
abstract deadlock patterns, instead of concrete ones. 
We thus expect our algorithms to be much more scalable in practice.


\SubParagraph{\underline{Deadlock-detection capability}}
In total, both \seqc and $\SyncPDOffline$ reported 40 deadlocks.
\seqc misses a deadlock of size $5$ in \texttt{DiningPhil},
which is detected by $\SyncPDOffline$,
and $\SyncPDOffline$ misses a deadlock in \texttt{jigsaw} which is detected by \seqc.
As $\SyncPDOffline$ is complete for sync-preserving deadlocks, we conclude that there are no more such deadlocks in our dataset.
Overall, $\SyncPDOffline$ and \seqc miss only three deadlocks reported by \dirk. 
On closer inspection, we found that these deadlocks are not witnessed by correct reorderings, and require reasoning about event values.
On the other hand, \dirk struggles to analyze even moderately-sized benchmarks and times out in $3$ of them. %(timeout is 3h).
This results in \dirk failing to report 5 deadlocks after $9$ hours, all of which are reported by $\SyncPDOffline$ in under a minute.
Similar conclusions were recently made in~\cite{Cai2021}.  
Overall, our results strongly indicate that the notion of sync-preservation characterizes most deadlocks that other tools are able to predict.


\SubParagraph{\underline{Unsoundness of \dirk}}
In our evaluation, we discovered that the soundness guarantee 
underlying \dirk~\cite{Kalhauge2018} is broken, resulting in it reporting false positives.
% Although \dirk is portrayed as sound, we have found two independent sources of unsoundness 
% resulting in it reporting false positives.
First, its constraint formulation~\cite{Kalhauge2018} 
does not rule out deadlock patterns when acquire events in the pattern hold common locks, 
in which case mutual exclusion disallows such a pattern to be a real predictable deadlock.
Second, \dirk also models conditional statements, allowing it to reason about witnesses beyond correct reorderings.
While this relaxation allows \dirk to predict additional deadlocks in \texttt{Transfer}, \texttt{Deadlock} and \texttt{HashMap}, 
its formalization is not precise and its implementation is erroneous.
We elaborate these aspects further in\begin{pldi}~\cite{arxiv}. \end{pldi}\begin{arxiv}~\appref{unsound-dirk}.\end{arxiv}


%We noticed that \dirk reports false positives in certain cases, which contradicts its theoretical soundness guarantee.
%We provide two such cases in~\appref{unsound-dirk}.
%The first case is a modified version of \texttt{Transfer}.
%Here, \dirk falsely reports a deadlock because it inadequately models conditional statements, which are used to allow more trace reorderings that can expose a deadlock.
%This relaxation enables \dirk to predict deadlocks in the benchmarks \texttt{Transfer} and \texttt{Deadlock}, which are missed by $\SyncPDOffline$ and \seqc. 
%However, this relaxation is not formalized and its implementation is erroneous.
%In the second case, \dirk falsely reports a deadlock due to missing that a common lock protects an otherwise deadlock pattern.

\SubParagraph{\underline{Running time}}
Our experimental results indicate that \dirk, backed by SMT solving, 
is the least efficient technique in terms of running time ---
it takes considerably longer or times out on large benchmark instances.
$\SyncPDOffline$ analyzed the entire set of traces $\sim\!\!\!21\times$ faster than \seqc.
On the most demanding benchmarks, such as 
HashMap and TreeMap, $\SyncPDOffline$ is more than $200\times$ faster than \seqc.
Although \seqc employs a polynomial-time algorithm for deadlock prediction, 
and thus significantly faster than the SMT-based \dirk,
the large polynomial complexity in its running time hinders scalability on 
execution traces coming from benchmarks that are more representative of realistic workloads.
In contrast, the linear time guarantees of $\SyncPDOffline$ are realized in practice, 
allowing it to scale on even the most challenging inputs.
More importantly, the improved performance comes while preserving essentially the same precision.


%\begin{figure}[!h]
\def\scatterscale{0.192}
\def\scatterwidth{0.23\textwidth}

\centering
\begin{subfigure}[b]{\scatterwidth}
\includegraphics[scale=\scatterscale]{figures/LinkedListDeadlockTest.pdf}
\label{subfig:sca-one-lock-clique}
\caption{LinkedList}
\end{subfigure}
\begin{subfigure}[b]{\scatterwidth}
\includegraphics[scale=\scatterscale]{figures/StackDeadlockTest.pdf}
\caption{Stack}
\label{subfig:sca-skewed}
\end{subfigure}


\caption{
Scalability experiments.
}
\label{fig:scalability}
\end{figure}


%\Andreas{We might revisit/remove the following}

%\textcolor{red}{check the numbers again.}
\SubParagraph{\underline{False negatives}}
Our benchmark set contains $93$ abstract deadlock patterns, $40$ of which are confirmed sync-preserving deadlocks.
We inspected the remaining $53$ abstract patterns to see if any of them are predictable deadlocks
missed by our sync-preserving criterion, independently of the compared tools.
$48$ of these $53$ patterns are in fact not predictable deadlocks ---
for every such pattern $D$, 
the set $S_D$ of events in the downward-closure of $\prev{}(D)$ with respect to $\tho{}$ and $\rf{}$,
already contains an event from $D$, disallowing any correct reordering
(sync-preserving or not) in which $D$ can be enabled.
Of the remaining, $4$ deadlock patterns obey the following scheme:
there are two acquire events $\acq_1, \acq_2$ participating in the deadlock pattern, 
each $\acq_i$ is preceded by a critical section on a lock that appears in 
$\lheld{}(\acq_{3-i})$, again disallowing a correct reordering that witnesses the pattern.
Thus, \emph{only one} predictable deadlock is not sync-preserving in our whole dataset.
This analysis supports that the notion of sync-preservation is not overly conservative in practice.
%\hunkar{I think here we should clarify that this analysis is not overly conservative as far as predictable deadlocks go.}

%\hunkar{
The above analysis concerns false negatives wrt. predictable deadlocks.
Some deadlocks are beyond the common notion of predictability we have adopted here, as they can only be exposed by reasoning about event values and control-flow dependencies, a problem that is $\NP$-hard even for 3 threads~\cite{Gibbons1997}.
We noticed $3$ such deadlocks in our dataset, found by \dirk,
though, as mentioned above, \dirk's reasoning for capturing such deadlocks is unsound in practice.
%}

%\hunkar{Maybe stress again that this reasoning about event values is non-trivial as it relies on heavyweight SMT solving. Also, maybe have the section "Unsoundness of Dirk" after this one and say that next we show that it is also tricky to implement in practice.}
%sync-preserving deadlocks are likely to be permissive and not lead to a high false negative rate.
%We conduct an analyses based on our benchmark set and characterize potential false negatives of our technique.
%Recall the conditions imposed on correct reorderings (see \secref{prelim}).
%One of the main conditions imposed by the definition of a correct reordering is that
%every read event reads from the same write as in the original trace. 
%We investigated the effect of this condition on the potential false negatives.
%Our benchmark set contains $93$ abstract deadlock patterns and our technique is able to find $42$ deadlocks.
%The remaining $51$ deadlock patterns constitute potential false negatives.
%If we only impose the above condition, then $46$ of the deadlock patterns cannot be realized to a real deadlock.
%The additional conditions that are specific to sync-preserving deadlocks (e.g., the order of critical sections on the same lock cannot be reversed) prevents the remaining $5$ deadlock patterns from being realizable.
%We remark that the definition of correct reorderings adopted in this work originate from the standard model that is widely used in this domain~\cite{Smaragdakis12,serbanuta2013,Koushik05}.
%Hence, this analyses further supports that given this standard program model, the additional restrictions introduced with the
%sync-preserving deadlocks are likely to be permissive and not lead to a high false negative rate.


\subsection{Online Experiments}
\seclabel{online-expr}

\Paragraph{Experimental setup}
%Dynamic analyses can incur high runtime overheads.
The objective of our second set of experiments is to evaluate 
the performance 
of our proposed algorithms in an \emph{online} setting.
For this, we implemented our $\SyncPDOnline$ algorithm inside the
framework of \dlfuzzer~\cite{Joshi2009} following closely the pseudocode in \algoref{online}. 
This framework instruments a concurrent program so that it can
perform analysis on-the-fly while executing it.
If a deadlock occurs during execution, it is reported and the execution halts.
However, if a deadlock is predicted in an alternate interleaving, 
then this deadlock is reported and the execution continues to search further deadlocks.
We used the same dataset as in \secref{offline-expr}, 
after discarding some benchmarks that could not be instrumented by \dlfuzzer.

To the best of our knowledge, all prior deadlock prediction techniques work offline.
For this reason, we only compared our online tool with the randomized 
scheduling technique of~\cite{Joshi2009} already
implemented inside the same \dlfuzzer framework.	
At a high level, this random scheduling technique works as follows.
Initially, it
(i)~executes the input program with a random scheduler, 
(ii)~constructs a \emph{lock dependency relation}, and 
(iii)~runs a cycle detection algorithm to discover deadlock patterns. 
For each deadlock pattern thus found, it spawns new executions that attempt 
to realize it as an actual deadlock.
To increase the likelihood of hitting the deadlock,
\dlfuzzer biases the random scheduler by pausing threads at specific locations.
%(e.g., before acquiring a certain lock).

The second, confirmation phase of~\cite{Joshi2009}
acts as a best-effort proxy for sound deadlock prediction.
On the other hand, $\SyncPDOnline$ is already sound and predictive, and thus does not require
additional confirmation runs, making it more efficient.
Towards effective prediction, we also implemented a simple bias to the scheduler.
If a thread $t$ attempts to write on a shared variable $x$ while holding a lock, then 
our procedure randomly decides to pause this operation for a short duration.
This effectively explores race conditions in different orders.
Overall, implementing $\SyncPDOnline$ inside \dlfuzzer provided the added advantage of supplementing a powerful prediction technique with a biased randomized scheduler.
%As an added advantage of implementing our algorithms in the
% framework, we are able to achieve
%the complementary objective of evaluating how well prediction supplements
%the effectiveness of an otherwise naive
%controlled concurrency testing
%technique like randomized scheduling.
To our knowledge, our work is the first to effectively 
combine these two orthogonal techniques.
We also remark that such a bias is of no benefit to \dlfuzzer itself
since it does not employ any predictive reasoning.

% a randomized testing procedure, even though the latter is rather simple.

For this experiment, we run \dlfuzzer on each benchmark, and for each deadlock pattern found in the initial execution, 
we let it spawn $3$ new executions trying to realize the deadlock, 
as per standard (\href{https://github.com/ksen007/calfuzzer}{https://github.com/ksen007/calfuzzer}).
We repeated this process $50$ times and recorded the total time taken.
Then, we allocated the same time for $\SyncPDOnline$ to repeatedly execute the same program and perform deadlock prediction.
We measured all deadlocks found by each technique.
%We also noticed that \dlfuzzer fails to work properly if the given input program deadlocks in the initial run.
%This results in \dlfuzzer itself to go into deadlock without producing any deadlock reports.  
%Hence, we made a modest modification on \dlfuzzer allowing it to work properly and report deadlocks in such cases.
%We used this modified version of \dlfuzzer in our evaluation.





%!TEX root = ../main.tex

\begin{table*}
\caption{
Performance comparison of $\SyncPDOnline$ ($\syncpdshort$) and $\dlfuzzer$ ($\dlfshort$).
%%Column 1 denotes the name of the benchmark.
Columns 2-3 show the total number of bug reports. 
Columns 4-6 show the total number of unique bugs found by each tool, and their union. 
Columns 7-12 show the hit rate on each bug.
Columns 13-16 show the runtime overhead of the tools.
%We refer to $\dlfuzzer$ as $\dlfshort$, to $\SyncPDOnline$ as $\syncpdshort$, and the instrumentation phases using the suffix $\mathsf{\tt -I}$.
%Performance comparison of $\SyncPDOnline$ ($\syncpdshort$) and $\dlfuzzer$ ($\dlfshort$).
%Column 1 denotes the name of the benchmark.
%Columns 2-3 show the total number of bug reports. 
%Columns 4-6 show the total number of unique bugs found by each tool, and their union. 
%Columns 7-12 show the hit rate on each bug.
%Columns 13-17 show the runtime overhead of the tools.
%We use the suffix $\mathsf{\tt -I}$ for the instrumentation phases.
\label{tab:expr-dlf-results}
}
\vspace{-0.15cm}
\setlength\tabcolsep{3pt}
\renewcommand{\arraystretch}{0.9}
\small
\centering
\scalebox{1}{
\begin{tabular}{|r|c|c|c|c|c|c|c|c|c|c|c|c|c|c|c|c|c|c|c|c|c|c|c|c|c|c|c|}
\hline
1 & 2 & 3 & 4 & 5 & 6 & 7 & 8 & 9 & 10 & 11 & 12 & 13 & 14 & 15 & 16 \\
\hline
\multirow{2}{*}{\textbf{Benchmark}}& 
\multicolumn{2}{c|}{ {\textsf{Bug Hits}}} & 
\multicolumn{3}{c|}{ {\textsf{Unique Bugs}}} 
& \multicolumn{2}{c|}{ \textsf{Bug 1}} &
\multicolumn{2}{c|}{ \textsf{Bug 2}} &
\multicolumn{2}{c|}{ \textsf{Bug 3}} & \multicolumn{4}{c|}{ {\textsf{Runtime Overhead}}}  \\
\cline{7-16}
\cline{2-6}
& \textsf{$\syncpdshort$} & \textsf{$\dlfshort$} & \textsf{$\syncpdshort$} & \textsf{$\dlfshort$}  & All 
& \textsf{\textsf{$\syncpdshort$}} 
& \textbf{$\dlfshort$} 
& {\textbf{$\syncpdshort$}} 
&{\textbf{$\dlfshort$}} 
& {\textbf{$\syncpdshort$}}
&{\textbf{$\dlfshort$}} & {\textbf{$\syncpdshortinstr$}}  & {\textbf{$\syncpdshort$}}  & 
{\textbf{$\dlfshortinstr$}}  &
 {\textbf{$\dlfshort$}} \\
\hline
%\multirow{2}{*}{\textbf{Benchmark}}& \multirow{2}{*}{$\mathcal{N}$} & \multirow{2}{*}{$\mathcal{T}$} & \multirow{2}{*}{$\mathcal{M}$}
%& \multirow{2}{*}{$\mathcal{L}$} &
%\multirow{2}{*}{$\NumAcquires+\mathcal{R}$} &\multicolumn{2}{c||}{ {\seqc}} & \multicolumn{2}{c||}{ {\textsf{$\SyncPDOffline$}}} & \multicolumn{2}{c|}{ {\textsf{$\SyncPDOnline$}}} \\
%\cline{7-12}
%& & & & & &  \textbf{$\#$Deadlocks} & { \textbf{Time (s)}} &{ \textbf{$\#$Deadlocks}} & { \textbf{Time (s)}} &{ \textbf{$\#$Deadlocks}} &{ \textbf{Time (s)}} \\
\hline
Deadlock & 50 & 50 & 1 & 1 & 1 & 50 & 50 & - & - & - & - & 2$\times$ & 3$\times$ & 2$\times$ & 4$\times$\\
Picklock & 227 & 97 & 2 & 1 & 2 & 226 & 97 & 1 & 0 & - & - & 2$\times$ & 2$\times$ & 2$\times$ & 5$\times$\\
Bensalem & 355 & 32 & 2 & 1 & 2 & 8 & 0 & 347 & 32 & - & - & 2$\times$ & 2$\times$ & 2$\times$ & 6$\times$\\
Transfer & 54 & 50 & 1 & 1 & 1 & 54 & 50 & - & - & - & - & 2$\times$ & 2$\times$ & 1$\times$ & 4$\times$\\
Test-Dimmunix & 702 & 0 & 2 & 0 & 2 & 351 & 0 & 351 & 0 & - & - & 2$\times$ & 2$\times$ & 2$\times$ & 4$\times$\\
StringBuffer & 153 & 131 & 2 & 2 & 2 & 128 & 118 & 25 & 13 & - & - & - & - & - & -\\
Test-Calfuzzer & 177 & 44 & 1 & 1 & 1 & 177 & 44 & - & - & - & - & 2$\times$ & 2$\times$ & 2$\times$ & 4$\times$\\
DiningPhil & 162 & 100 & 1 & 1 & 1 & 162 & 100 & - & - & - & - & - & - & - & -\\
HashTable & 169 & 120 & 2 & 2 & 2 & 82 & 21 & 87 & 99 & - & - & - & - & - & -\\
Account & 19 & 188 & 1 & 1 & 1 & 19 & 188 & - & - & - & - & 2$\times$ & 8$\times$ & 2$\times$ & 16$\times$\\
Log4j2 & 290 & 100 & 2 & 1 & 2 & 145 & 100 & 145 & 0 & - & - & - & - & - & -\\
Dbcp1 & 265 & 138 & 2 & 2 & 2 & 264 & 61 & 1 & 77 & - & - & - & - & - & -\\
Dbcp2 & 129 & 126 & 2 & 2 & 2 & 86 & 99 & 43 & 27 & - & - & - & - & - & -\\
RayTracer & 0 & 0 & 0 & 0 & 0 & - & - & - & - & - & - & 122$\times$ & 124$\times$ & 109$\times$ & 111$\times$\\
Tsp & 0 & 0 & 0 & 0 & 0 & - & - & - & - & - & - & 47$\times$ & 60$\times$ & 37$\times$ & 40$\times$\\
jigsaw & 1189 & 1 & 1 & 1 & 2 & 1189 & 0 & 0 & 1 & - & - & - & - & - & -\\
elevator & 0 & 0 & 0 & 0 & 0 & - & - & - & - & - & - & 2$\times$ & 2$\times$ & 2$\times$ & 2$\times$\\
JDBCMySQL-1 & 349 & 117 & 2 & 3 & 3 & 1 & 21 & 0 & 4 & 348 & 92 & 3$\times$ & 4$\times$ & 2$\times$ & 13$\times$\\
JDBCMySQL-2 & 559 & 73 & 1 & 1 & 1 & 559 & 73 & - & - & - & - & 2$\times$ & 4$\times$ & 2$\times$ & 18$\times$\\
JDBCMySQL-3 & 560 & 224 & 1 & 1 & 1 & 560 & 224 & - & - & - & - & 2$\times$ & 5$\times$ & 2$\times$ & 24$\times$\\
JDBCMySQL-4 & 1717 & 101 & 3 & 1 & 3 & 95 & 0 & 834 & 0 & 788 & 101 & 3$\times$ & 5$\times$ & 2$\times$ & 31$\times$\\
hedc & 0 & 0 & 0 & 0 & 0 & - & - & - & - & - & - & 2$\times$ & 2$\times$ & 1$\times$ & 2$\times$\\
cache4j & 0 & 0 & 0 & 0 & 0 & - & - & - & - & - & - & 2$\times$ & 2$\times$ & 2$\times$ & 2$\times$\\
lusearch & 0 & 0 & 0 & 0 & 0 & - & - & - & - & - & - & 16$\times$ & 17$\times$ & 13$\times$ & 16$\times$\\
ArrayList & 47 & 45 & 3 & 3 & 3 & 20 & 22 & 3 & 5 & 24 & 18 & 50$\times$ & 69$\times$ & 18$\times$ & 79$\times$\\
Stack & 44 & 27 & 3 & 3 & 3 & 18 & 13 & 8 & 4 & 18 & 10 & 69$\times$ & 91$\times$ & 64$\times$ & 86$\times$\\
IdentityHashMap & 68 & 62 & 2 & 2 & 2 & 13 & 47 & 55 & 15 & - & - & 4$\times$ & 8$\times$ & 3$\times$ & 10$\times$\\
LinkedList & 48 & 26 & 3 & 2 & 3 & 21 & 17 & 7 & 0 & 20 & 9 & 16$\times$ & 28$\times$ & 14$\times$ & 32$\times$\\
Swing & 0 & 0 & 0 & 0 & 0 & - & - & - & - & - & - & 5$\times$ & 6$\times$ & 4$\times$ & 6$\times$\\
Sor & 0 & 0 & 0 & 0 & 0 & - & - & - & - & - & - & 2$\times$ & 7$\times$ & 2$\times$ & 2$\times$\\
HashMap & 46 & 44 & 2 & 2 & 2 & 18 & 11 & 28 & 33 & - & - & 7$\times$ & 11$\times$ & 4$\times$ & 8$\times$\\
Vector & 126 & 50 & 1 & 1 & 1 & 126 & 50 & - & - & - & - & 2$\times$ & 2$\times$ & 2$\times$ & 3$\times$\\
LinkedHashMap & 57 & 43 & 2 & 2 & 2 & 22 & 10 & 35 & 33 & - & - & 10$\times$ & 10$\times$ & 4$\times$ & 8$\times$\\
WeakHashMap & 29 & 40 & 2 & 2 & 2 & 6 & 11 & 23 & 29 & - & - & 7$\times$ & 12$\times$ & 4$\times$ & 8$\times$\\
montecarlo & 0 & 0 & 0 & 0 & 0 & - & - & - & - & - & - & 16$\times$ & 100$\times$ & 13$\times$ & 126$\times$\\
TreeMap & 42 & 47 & 2 & 2 & 2 & 16 & 15 & 26 & 32 & - & - & 9$\times$ & 12$\times$ & 5$\times$ & 9$\times$\\
eclipse & 0 & 0 & 0 & 0 & 0 & - & - & - & - & - & - & 2$\times$ & 2$\times$ & 2$\times$ & 2$\times$\\
TestPerf & 0 & 0 & 0 & 0 & 0 & - & - & - & - & - & - & 2$\times$ & 2$\times$ & 2$\times$ & 2$\times$\\
\hline\hline\textbf{Total} & \textbf{7633} & \textbf{2076} & \textbf{49} & \textbf{42} & \textbf{51} & - & - & - & - & - & - & - & - & - & - \\ 
\hline
\end{tabular}
}
\end{table*}


\Paragraph{Evaluation}
\cref{tab:expr-dlf-results} presents our experimental results.
%A bug identifies a unique tuple of source code locations corresponding to events
%participating in the deadlock.
Columns $2$-$3$ of the table display the total number of bug hits,
which is the total number of times a bug was predicted by $\SyncPDOnline$ in the entire duration,
or was confirmed in any trial of \dlfuzzer.
Columns $4$-$6$ display the unique bugs (i.e., unique tuples of source code locations leading to a deadlock) 
found by the techniques.
The employed techniques are able to find a maximum of $3$ unique bugs for each benchmark
in our benchmark set. 
Respectively, columns $7$-$12$ display the detailed information on the number 
of times a particular bug was found by each technique.
Runtime overheads are displayed in the columns $13$-$16$, with $\mathsf{\tt -I}$ denoting the instrumentation phase only.


\SubParagraph{\underline{Deadlock-detection capability}}
\dlfuzzer had $2076$ bug reports in total, and it found $42$ unique bugs.
In contrast, $\SyncPDOnline$ flagged $7633$ bug reports, corresponding to $49$ unique bugs.
In more detail, \dlfuzzer missed $9$ bugs reported by \SyncPDOnline whereas 
$\SyncPDOnline$ missed $2$ bugs reported by \dlfuzzer.
Also, \SyncPDOnline significantly outperformed \dlfuzzer in total number of bugs hits.
Our experiments again support that the notion of sync-preservation  captures most deadlocks that occur in practice, to the extent that other state-of-the-art techniques can capture.
%\hunkar{
A further observation is that in the offline experiments, \SyncPDOffline  was not able to find deadlocks in \texttt{Transfer} and \texttt{Deadlock}. 
However, the random scheduling procedure allowed \SyncPDOnline 
to navigate to executions from which deadlocks can be predicted.
This demonstrates the potential of combining predictive dynamic
techniques with controlled concurrency testing.
%; a direction we find promising to be pursued further by the community.


\SubParagraph{\underline{Runtime overhead}}
We have also measured the runtime overhead of both \SyncPDOnline and \dlfuzzer,
both as incurred by instrumentation, as well as by the deadlock analysis.
The latter is the time taken by \algoref{online} for the case of $\SyncPDOnline$,
and the overhead introduced due to the new executions in the second confirmation phase for the case of \dlfuzzer.
Our results show that the instrumentation overhead of \SyncPDOnline is, in fact, comparable to that of \dlfuzzer, though somewhat larger. 
This is expected, as \SyncPDOnline needs to also instrument memory access events, while \dlfuzzer only instruments lock events, but at the same time surprising because the number of memory access events
is typically much larger than the number of lock events.
On the other hand, the analysis overhead is often larger for \dlfuzzer, 
even though it reports fewer bugs.
It was not possible to measure the runtime overhead in certain benchmarks as 
either they were always deadlocking or the computation was running indefinitely.





\section{Conclusion and Future Work} 
\section{Conclusions}
In this paper, we set out to address the problem of multi-tasking robots in multi-robot tasks. 
%A fundamental limitation of existing multi-robot systems was addressed by the removal of a restrictive assumption that was often made--robots are single-tasking.
%Our method allowed coalitions to overlap thus enabling multi-tasking robots. 
We observed that the key underlying challenge was to reason about the physical constraints that could be synergistically satisfied.
%which directly affected the feasibility of multi-tasking.
To address the challenge, we developed our method based on the information invariant theory and modeled constraints as information instances. 
%This allowed us to reason about the relationships between constraints by reasoning about those between information requirements. 
Thereby, a formal and general framework to achieve multi-tasking robots was developed. 
We showed that our algorithm was sound and complete under our problem settings. 
%Our method was integrated with a simple greedy heuristic for task allocation.
Simulation  results  were  provided  to  show  the  effectiveness  of  our approach under resource-constrained situations and in handling challenging situations. % in a multi-UAV simulator. 

% The idea of multi-tasking is attractive in many ways. 
% Humans are living in multi-tasking environments--at any point of time, 
% we are optimizing for more than one task. 
% Multi-task often leads to more efficient task performance since it allows us to exploit task synergies. 
% The work presented in this paper takes us one step forward in realizing multi-tasking robots. 
% In particular, we started looking at the feasibility of multi-tasking. 
% There are many potential directions to pursue along this direction. First, several limitations are present in the current approach. 
% For example, although our method guarantees that there exists a physical configuration that satisfies all the constraints, it does not explicitly take the environmental influence into account. For example, a narrow corridor may prevent a robot formation from passing through, even though all the constraints for the formation do not introduce any conflicts. In this sense, our work should better be characterized as establishing a necessary condition for multi-tasking. Also, our method is mainly focused on the ``{\it planning}'' phase and hence does not address how the robots reach the desired configuration and maintain the constraints. These issues are assumed to be handled by the execution layer.

% More generally, the question of how to execute the tasks with overlapping coalitions is not addressed in this work. 
% As we already discussed, executing individual tasks with non-overlapping coalitions is straightforward but task synergies impose additional requirements on the task execution: how should the robots that are assigned multiple tasks execute them? Should they consider them in a prioritized strategy~\cite{van2005prioritized}? Or should they combine the different tasks in a way that is similar to motor schemas~\cite{arkin2}. 
% Communication requirements for maintaining the constraints must also be taken into account. How should the robots optimize their communication to improve the task performance? 

% The stringency of the physical constraints is another interesting question. It may be desirable to relax the constraints in certain situations (e.g., due to environmental influences). In such cases, it may be important to consider the problem where the constraints are least violated~\cite{kim2012revision}, or specify task constraints in different ways to increase the diversity of the configurations~\cite{srivastava2007domain} so as to make it robust to different environments. 






\balance %

\bibliographystyle{IEEEtran}
\bibliography{Bibs}

\appendix \label{sec:Appendix}
%% The Appendices part is started with the command \appendix;
%% appendix sections are then done as normal sections
\appendix
% \onecolumn
\subsection{Prompt Template}
\label{app1}

Figs. \ref{appendix:prompt1}-\ref{appendix:prompt3} (due to the page length, we break it into three parts) show the prompt design for the information extraction in the context of construction project scheduling. 

\begin{figure}[t]
    \vspace{-4.5cm}
    \centering
    \begin{tcolorbox}[colback=gray!10!white, colframe=gray!50!gray, halign=left, boxrule=0.5pt, left=1mm, right=1mm, top=1mm, bottom=1mm]
    \fontsize{8pt}{8pt}\selectfont
    SYSTEM PROMPT: You are a project management assistant specializing in construction scheduling analysis. Your task is to analyze text descriptions of project changes and extract structured information about task relation changes in a construction project.
    \vspace{8pt}
    
    CONTEXT\\
    The project involves the following tasks and their relationships: \\
    \vspace{3pt}
    Task ID \textbar{} Predecessor \textbar{} Duration \textbar{} Description \textbar{} Robot Type
    \begin{itemize}
    \item T1 \textbar{} - \textbar{} 0.25 \textbar{} Move Electrical Conduit \textbar{} R1
    \item T2 \textbar{} - \textbar{} 0.25 \textbar{} Move Window Frame \textbar{} R1
    \item T3 \textbar{} - \textbar{} 0.25 \textbar{} Move Window \textbar{} R1
    \item T4 \textbar{} - \textbar{} 0.25 \textbar{} Move Duct Structural Materials \textbar{} R1
    \item T5 \textbar{} - \textbar{} 0.25 \textbar{} Move Duct \textbar{} R1
    \item T6 \textbar{} - \textbar{} 0.5 \textbar{} Drill Wall \textbar{} R4 or R2
    \item T7 \textbar{} T1, T6 \textbar{} 1 \textbar{} Install Electrical Conduit \textbar{} R5 or R2
    \item T8 \textbar{} T2 \textbar{} 1 \textbar{} Install Window Frame \textbar{} R4 or R2
    \item T9 \textbar{} T3, T8 \textbar{} 0.5 \textbar{} Install Window \textbar{} R3
    \item T10 \textbar{} T4 \textbar{} 2 \textbar{} Duct Structural Framing \textbar{} R4 or R2
    \item T11 \textbar{} T5, T10 \textbar{} 2 \textbar{} Install HVAC Duct \textbar{} R4 or R2
    \item T12 \textbar{} T7 \textbar{} 2 \textbar{} Install Wiring \textbar{} R5 or R2
    \item T13 \textbar{} T12 \textbar{} 1 \textbar{} Wall Painting \textbar{} R6
    \item T14 \textbar{} - \textbar{} 0.5 \textbar{} Construction Site Inspection \textbar{} R7
    \end{itemize}
    \vspace{8pt}
    
    The robot capabilities are listed below: \\
    \vspace{3pt}
    Robot ID \textbar{} Capabilities
    \begin{itemize}
    \item R1: Cargo container
    \item R2: High-payload, Precise parallel gripper, Normal parallel gripper
    \item R3: High-payload, Suction-based gripper
    \item R4: High-payload, Normal parallel gripper
    \item R5: Precise parallel gripper
    \item R6: Sprayer
    \item R7: Camera, IAQ sensors
    \end{itemize}
    \vspace{8pt}
    
    CONSTRAINT TYPES:
    \begin{enumerate}
    \item Task Dependency Adjustments
      \begin{itemize}
        \item Format: [task\_id, successor, +/-]
        \item task\_id: the target task
        \item successor: the successors of the target task
        \item +/-: ``+'' indicates a newly added successor, ``-'' means the dependency has been removed
      \end{itemize}
    \item Task Duration Variations
      \begin{itemize}
        \item Format: [task\_id, new\_duration]
        \item task\_id: the target task
        \item new\_duration: the new duration of the target task in hours
      \end{itemize}
    \item Task Starting Time Changes
      \begin{itemize}
        \item Format: [task\_id, start\_time\_change]
        \item task\_id: the target task
        \item start\_time\_change: the changes in start time of the target task (e.g., +2 means delayed by 2 hours; -2 means ahead by 2 hours)
      \end{itemize}
    \item Number of Robot Variations
      \begin{itemize}
        \item Format: [robot\_type\_id, robot\_number\_change]
        \item robot\_type\_id: the type of robot (e.g., R1, R2, R3, etc.)
        \item robot\_number\_change: the number changes of the robot (e.g., +1 means one more robot; -1 means one less robot)
      \end{itemize}
    \item Task Conflict Constraints
      \begin{itemize}
          \item Format: [task\_id1, task\_id2]
          \item task\_id1: the first task in the conflict
          \item task\_id2: the second task in the conflict
      \end{itemize}
    \end{enumerate}
    \end{tcolorbox}
    \caption{Prompt design for construction project scheduling - Part 1.}
    \label{appendix:prompt1}
\end{figure}


\begin{figure}[t]
    \centering
    \begin{tcolorbox}[colback=gray!10!white, colframe=gray!50!gray, halign=left, boxrule=0.5pt, left=1mm, right=1mm, top=1mm, bottom=1mm]
    \fontsize{8pt}{8pt}\selectfont
    STEP-BY-STEP INSTRUCTIONS:
    \begin{enumerate}
    \item Read through the entire description to understand the context.
    \item For each change mentioned in the description:
       \begin{itemize}
       \item[a.] Identify which task (T1-T14) or robot type (R1-R7) is being affected based on CONTEXT. 
        \begin{itemize}
          \item Be careful to distinguish between similar tasks, for example:
          \begin{itemize}
          \item T2 (Move Window Frame) vs. T3 (Move Window) vs. T8 (Install Window Frame) vs. T9 (Install Window) - These are different tasks.
          \item If text mentions ``window installation'', specifically, it refers to T9 (Install Window), not T3 or T8
          \item If text mentions "window frame installation," it refers to T8 (Install Window Frame), not T2
          \end{itemize}
        \item Be careful to distinguish between similar robots, for example:
          \begin{itemize}
          \item R2, R3, R4, and R5 are different robots. 
          \item Only R2 combines both high-payload and precise parallel gripper capabilities.
          \item If text only mentions ``high-payload and normal parallel gripper'', it refers to R4 not R2.
          \end{itemize}
       \end{itemize}
       \item[b.] Determine which constraint type (1-5) applies to the change based on CONSTRAINT TYPES.
       \item[c.] Extract the specific parameters needed for that constraint type.
       \item[d.] Format the parameters according to the required format for that constraint type.
       \end{itemize}
    \item Compile all identified changes into the JSON output format:
       \begin{itemize}
       \item[a.] Create a JSON object with a ``changes'' array.
       \item[b.] For each change, add an object with ``constraint\_type'' and ``parameters'' fields.
       \item[c.] Ensure numerical values (like durations and time changes) are formatted as numbers, not strings.
       \item[d.] Ensure task IDs, successors, and robot types are formatted as strings.
       \item[e.] For time-related values:
          \begin{itemize}
          \item Simplify all numerical values to their simplest form (e.g., 1.5 not 1.50, 2 not 2.0)
          \item Convert minutes to hours (e.g., 30 minutes = 0.5 hours, 45 minutes = 0.75 hours)
          \item Please be aware that if you identify the constraint as 3, the time change should be associated with ``+'' or ``-''. 
          \end{itemize}
       \item[g.] Please be aware that if you identify the constraint as 4, the robot change should be associated with ``+'' or ``-''.
       \end{itemize}
    \item Double-check your result to ensure all changes mentioned in the description have been captured.
       \begin{itemize}
       \item[a.] Please ensure that your output follows the required format; e.g., for constraint 1, the output should be [task\_id, successor, +/-] (do NOT nest successors in additional brackets) and the task\_id should be the predecessor of the successor. 
       \item[b.] Please ensure that if you identify the constraint as 1, you correctly identify the target task and the successor of the target task and put them in the right order [task\_id, successor, +/-].
       \item[c.] Please ensure that if you identify the constraint as 3, the time change should be associated with ``+'' or ``-''. 
       \item[d.] Please ensure that if you identify the constraint as 4, the robot change should be associated with ``+'' or ``-''. 
       \item[e.] Please ensure that the task description corresponds to the task\_id in the CONTEXT.
       \end{itemize}
    \end{enumerate}
    \end{tcolorbox}
    \caption{Prompt design for construction project scheduling - Part 2.}
    \label{appendix:prompt2}
\end{figure}


\begin{figure}[t]
    \vspace{-2cm}
    \centering
    \begin{tcolorbox}[colback=gray!10!white, colframe=gray!50!gray, halign=left, boxrule=0.5pt, left=1mm, right=1mm, top=1mm, bottom=1mm]
    \fontsize{8pt}{8pt}\selectfont
EXAMPLES:\\
    Example 1:\\
    Input: ``Due to how things are unfolding on-site, it's understood that the drilling machine is not functioning, so the wall will be drilled manually. The task is expected to take two hours, and in light of recent discussions, after coordinating with field staff, it seems that the original worker assigned to install the HVAC duct is no longer available; however, we have secured another worker who can arrive in 150 minutes.''\\
    Output:
    \begin{verbatim}
{"changes": [
{"constraint_type": 2, "parameters": [T6, 2]},
{"constraint_type": 3, "parameters": [T11, +2.5]}
]}
    \end{verbatim}

    Example 2:\\
    Input: ``Recent developments suggest that wall painting takes 1.5 hours instead of 1 hour due to the need for multiple coats, and in light of recent adjustments, a revised understanding across teams indicates that a specialist required for electrical conduit installation calls in sick, preventing work from starting for 2 hours., followed by further refinements as recent developments suggest that wall painting takes 1.5 hours instead of 1 hour due to the need for multiple coats.''\\
    Output:
    \begin{verbatim}
{"changes": [
{"constraint_type": 2, "parameters": [T13, 1.5]},
{"constraint_type": 3, "parameters": [T7, +2]},
{"constraint_type": 2, "parameters": [T13, 1.5]}
]}
    \end{verbatim}

    Example 3:\\
    Input: ``Task dependencies have shifted, and one of the robots capable of handling heavy loads and performing fine, precise tasks is currently out of service due to a mechanical failure. Additionally, in light of recent discussions and the evolving situation on-site, it appears that two robots with high-capacity arms and fine-movement grippers were not charged, and have now run out of power.''\\
    Output:
    \begin{verbatim}
{"changes": [
{"constraint_type": 4, "parameters": [R2, -1]},
{"constraint_type": 4, "parameters": [R2, -2]}
]}
    \end{verbatim}
    
    Now, analyze the following description and extract all task relation changes in the specified JSON format: \{description\}
    
    Please output your response in JSON format and do not output other things. 
    \begin{verbatim}
{"changes": [
{"constraint_type": <number>, 
 "parameters": [<value1>, <value2>, ...]},...
]}
    \end{verbatim}
    \end{tcolorbox}
    \caption{Prompt design for construction project scheduling - Part 3.}
    \label{appendix:prompt3}
\end{figure}

\end{document}
